
    \documentclass[a4paper]{article}
    \usepackage{savetrees}
    \usepackage{hyperref}
    \usepackage{pdfpages}

    \begin{document}

    Ryan Wang, Matthew Finlayson, Luca Soldaini, Swabha Swayamdipta, and Robin Jia.  {\bf \href{http://cs.umd.edu/~jbg//docs/2025_colm_slung.pdf}{Teaching Models to Understand (but not Generate) High-risk Data}}.  \emph{Conference on Language Modeling}, 2025.



    \vspace{5cm}

    \begin{verbatim}
    @inproceedings{Swayamdipta,:Jia-2025,
	Title = {Teaching Models to Understand (but not Generate) High-risk Data},
	Author = {Ryan Wang, Matthew Finlayson, Luca Soldaini, Swabha Swayamdipta, and Robin Jia},
	Booktitle = {Conference on Language Modeling},
	Year = {2025},
	Url = {http://cs.umd.edu/~jbg//docs/2025_colm_slung.pdf},
}

    \end{verbatim}

    {\bf Accessible Abstract:} Language models are pre-trained on large text corpora which are filtered to remove undesired content that we do not want the models to generate. However, filtering can come with a cost: models will not learn to understand the topics that are filtered out. We propose a new method, SLUNG, which enables models to learn to understand some pieces of training data without learning to emulate it during generation. This enables models to acquire more diverse, wide-ranging knowledge while still controlling its behavior; for example, language models can learn to detect when a statement is toxic, but never generate toxic speech itself.\\~

    

    

    \vspace{3cm}
    Downloaded from \url{http://cs.umd.edu/~jbg/docs/2025_colm_slung.pdf}  \\

    \vspace{2cm}
    \textit{Contact Jordan Boyd-Graber (jbg@boydgraber.org) for questions about this paper.}

    \includepdf[pages={-}]{src_docs/2025_colm_slung.pdf}

    \end{document}
    